\chapter{Implementation and Test Cases}
\label{ch:implementation}
The implementation of FYP-I solves the problem of selection of a language model in AetherFlow with the consideration of factors such as quality, cost, latency, and privacy that influence the selection process. The AetherFlow routing prototype is available at \url{https://aether-flow-gamma.vercel.app/}. The implementation includes a web interface, JSON routing API, and some preliminary benchmarking results. The provided benchmarking results have not been independently reproduced and cannot be considered verified performance results. The other part of FYP-II includes VS Code integration, semantic caching, temporary value protection, and usage logging.

\section{Implementation}

\subsection{Web Interface and API}
The UI is constructed using React and loads the model list, routing configuration, and sample queries. The user inputs the query, task type, and mode. The request is sent to \texttt{/api/route} endpoint by the UI. It then presents the selected model, prediction of answer quality, cost estimation, response-time information, and the rationale for selecting the model.

The API is built using FastAPI and comes with interactive documentation and an OpenAPI specification of its input/output parameters. The routing request includes the query, along with optional \texttt{query\_class}, \texttt{router}, \texttt{mode}, \texttt{threshold}, \texttt{quality\_floor}, \texttt{latency\_budget\_ms}, and \texttt{sensitive}. These parameters define the query, routing configuration, and privacy option. Input validation logic will be documented and tested independently of the gateway integration process.

\begin{table}[!htbp]
\centering\small
\caption{Public interfaces of the implemented routing demonstrator}
\renewcommand{\arraystretch}{1.20}
\setlength{\extrarowheight}{1pt}
\begin{tabular}{|P{0.30\textwidth}|P{0.60\textwidth}|}
\hline
\textbf{Endpoint} & \textbf{Role} \\ \hline
\texttt{POST /api/route} & Choose a model using Legacy routing or a selected policy mode. \\ \hline
\texttt{POST /api/sensitivity} & Check whether a query matches the current sensitive-data rules. \\ \hline
\texttt{GET /api/models} & Return the model list and reported benchmark values, which still need checking. \\ \hline
\texttt{GET /api/results} & Return early comparison results and dataset details that still need checking. \\ \hline
\texttt{GET /api/objectives} & Return policy settings and evaluation summaries that still need checking. \\ \hline
\texttt{GET /api/pareto} & Show reported model choices with different quality, cost, and speed trade-offs. The measurements still need checking. \\ \hline
\texttt{GET /api/privacy} & Return privacy settings and the models allowed for sensitive queries. \\ \hline
\end{tabular}
\label{tab:implemented_api}
\end{table}

\subsection{Model List and Query Information}
The saved model-list response dated 11 November 2026 consists of eight models: Llama-3.1-8B-Instruct, Qwen3-8B, DeepSeek-v3-0324, Gemini-2.5-Flash, GPT-4.1, Claude-Sonnet-4, Gemini-2.5-Pro, and GPT-5. The name of a model on the list means that the model is set up in the prototype. The name is no guarantee of whether the model is available and connected to respond to queries. The list of models also includes cost, response time, and accuracy metrics for each task type. These metrics have not been independently verified.

There are five task types used in the prototype: Coding, Competitive Math, General Knowledge, Mathematical Reasoning, and Scientific Questionnaire. The routing response includes estimated task type, probability of the query being easy, medium, or hard, and probability of getting a correct response for each model. This information helps in selecting a model prior to sending the query to a model. Information about training and scoring adjustment procedure are included in the API documentation; however, the source code, the weights of the trained model, the dataset, and the logs have not been independently verified.

\subsection{Legacy Model Selection}
Legacy routing compares each model's predicted chance of answering correctly with a configured confidence requirement. If no candidate meets that requirement, the service returns its configured fallback model. The response includes the selected model, the fallback flag, predicted scores, and reasons for the decision. These fields describe the prototype's selection output; they do not establish that the selected model generated a correct answer. The preliminary observations are discussed in Chapter~\ref{ch:results}.

\subsection{Routing Modes}
The second router can handle Economy, Balanced, Speed, Quality, and Private modes. The request picks the router with \texttt{router=multi\_objective}, along with a necessary mode. Economy mode pays attention to cost more. Balanced mode takes into account both cost and speed. Speed mode uses an estimate of the maximum response time. Quality mode is oriented at predicted correctness. The Private mode permits only those models which meet the privacy configuration requirements.

The configurations include weights for cost and response time significance. They also include scales and adjustments for measuring values in order to put various measures on the same scale and modify predicted scores. The way these values were obtained is not independently verified. Setting a quality objective does not ensure that the chosen model will have such level of accuracy. Modes described at \texttt{/api/modes} differ from the five modes at \texttt{/api/objectives}.

\subsection{Privacy Rules and Sensitive Queries}
The current configuration permits only Llama-3.1-8B-Instruct and Qwen3-8B for queries that are deemed to be sensitive. In the preliminary experiments, in both Private mode and Balanced with the query being sensitive, an allowed model was chosen. This holds for the tested requests; it does not mean that all queries or fallbacks will comply with the policy. There are six patterns and keywords in the filtering endpoint configuration. Some of the keywords include confidential, medical record, and proprietary. These guidelines offer a basic filtering process, but the detection rate of such policies has not been quantified.

The configuration may also prohibit sensitive queries and/or turn off the use of external models. Provider data retention policy is provided by the administrator and cannot confirm the actual storage of the data by the provider. Choosing an appropriate model is just a routing decision. Further analysis is required to confirm local execution of the model, no private data leakage, removal of the temporary mapping, and no logging of sensitive data.

\subsection{Returned Results and Remaining Integration}
A multi-objective routing response consists of the selected model, mode, estimated quality, routing score, estimated cost, and response-time information. It distinguishes the decision time of the router from the estimated response time of the selected model. The latter uses average values from the benchmark, not a response produced by a model in response to the request. The measurement technique and values used in the measurement have not been independently validated. Thus, the resulting value is an estimate and not a verified response time. The model name produced by the router does not indicate whether the provider was actually called.

Now the task is to integrate the selector with the FastAPI gateway and VS Code client. The gateway will validate models' authorization, retrieve a cached response if any, select a model if necessary, invoke it, and return the response. The use of PostgreSQL for logging, Qdrant for caching, the Privacy Vault, authentication of the user, and error handling are yet to be implemented.

\subsection{Request Processing Steps}
The prototype produces the routing decision for the provided query and parameters. The saved experimental responses identify the selected model and related decision fields for the supplied inputs. For Legacy routing, the service estimates model scores based on the task type and difficulty. Then, the service compares the scores with the threshold and picks a model or the default fallback one. For multi-objective routing, the service takes into account the mode and privacy, discards the inappropriate models, compares the valid candidates based on the mode parameters and gives the selected model and explanation. This is the description of the functionality that is revealed through the API; it is not validated at all levels of internal training or scoring implementations.

Privacy criteria are the requirements for the model, while cost and speed weights are preferences for comparing the allowed models. An inexpensive external model should not become allowed due to high scores. The Speed mode additionally considers an estimate of the acceptable response time. The preliminary experiments examine particular routing decisions only. There are no broader tests of failures, fallbacks and full gateway yet.

\section{Conclusion}
The FYP-I implementation provides a web interface and an API for exploring model-selection decisions. The available responses expose the configured model list, routing choices, privacy settings, and simple keyword-screening explanations. These components provide a starting point for the proposed gateway, while the training records and reported performance summaries still require independent verification. Chapter~\ref{ch:results} examines the preliminary experimental evidence. Live inference, the semantic cache, the Privacy Vault, usage logging, and VS Code integration remain work for FYP-II, together with formal integration and acceptance testing.

\clearpage


\chapter{Experimental Results and Discussion}
\label{ch:results}

This chapter describes the initial experiments carried out using the present AetherFlow routing prototype. The experiments explore issues of configuration of the model, influence of routing parameters on the selected model, and reaction of the keyword screening interface. The findings are based on the analysis of recorded API request-response pairs and will be evaluated with respect to expected functioning of the prototype. The focus of this evaluation is selection of models and interface output. The performance of answer generation and the whole gateway are not evaluated here.

\section{Experimental Setup}
The prototype was assessed by providing selected prompts with various routing configurations and analyzing the resulting output. The tests concentrated on model selection and sensitivity detection. The prompt ``Explain binary search.' was provided with the use of Legacy, Balanced, and Private modes. The artificial prompt ``This is confidential proprietary information.' was provided with the help of Balanced mode and with the parameter \texttt{sensitive=true}. This prompt was also used with the sensitivity endpoint. The saved API requests were passed via the corresponding Python script with one-time usage of every input and configuration. There are seven responses analyzed here: one model list, one privacy configuration, four routing decisions, and one sensitivity detection. No actual personal data was provided.

\section{Preliminary Experimental Results}
The experiments try to answer the following questions: if the prototype reveals the configuration required for the selection, if the settings of the routes influence the choice returned, and if the configuration of a sensitivity keyword results in the expected classification. Seven out of seven applicable requests have resulted in HTTP 200. This proves that the interfaces returned positive responses to these requests; the interpretation of their content will be done below.

\subsection{Model Configuration and Approved Destinations}
Experiment 1 retrieved the list of models and the privacy parameters. The response for model list had eight unique, non-empty model IDs. The response for privacy parameters returned two models that are approved for sensitive queries. Table~\ref{tab:prototype_configuration_results} shows the observed configuration.

\begin{table}[!htbp]
\centering\small
\setstretch{1}\renewcommand{\arraystretch}{1.15}
\caption{Observed model and privacy configuration}
\begin{tabular}{|P{0.23\textwidth}|P{0.30\textwidth}|P{0.35\textwidth}|}\hline
\textbf{Configuration} & \textbf{Recorded Result} & \textbf{Interpretation} \\ \hline
Model list & Eight distinct model identifiers were returned. & The interface exposes a defined set of model entries against which routing choices can be checked. \\ \hline
Sensitive-query approvals & Llama-3.1-8B-Instruct and Qwen3-8B were approved. & The configuration supplies an explicit allowed set for the privacy-related comparisons. \\ \hline
\end{tabular}
\label{tab:prototype_configuration_results}
\end{table}

These answers define the setup for interpreting the routing experiment. For example, a chosen identifier can be verified from the registry, and a privacy-oriented decision can be checked against the approved models. Nevertheless, a configured identifier does not confirm that a model is connected, accessible, and active in this setting. In the same way, the approval flag represents a configured rule and does not indicate anything about the provider’s handling of the data. Thus, the experiment shows a reference configuration rather than a benchmark of model capabilities or an example of its secure operation.

\subsection{Effect of Routing Settings on Model Selection}
The second experiment compared answers to the ordinary prompt under three routing configurations. Legacy routing resulted in GPT-5. Balanced routing resulted in GPT-4.1. Private mode resulted in Llama-3.1-8B-Instruct. Additional Balanced routing request labeled as sensitive resulted in Llama-3.1-8B-Instruct. Chosen identifiers were part of the recorded list of models. Table~\ref{tab:verified_routes} contains the choices without considering quality, cost, and latency estimates as measured results.

\begin{table}[!htbp]
\centering\small
\setstretch{1}\renewcommand{\arraystretch}{1.15}
\caption{Returned model choices under the recorded routing settings}
\begin{tabular}{|P{0.26\textwidth}|P{0.30\textwidth}|P{0.32\textwidth}|}\hline
\textbf{Prompt} & \textbf{Routing Setting} & \textbf{Selected Model} \\ \hline
Ordinary sample & Legacy & GPT-5 \\ \hline
Ordinary sample & Balanced & GPT-4.1 \\ \hline
Ordinary sample & Private & Llama-3.1-8B-Instruct \\ \hline
Synthetic sensitive sample & Balanced with sensitivity flag & Llama-3.1-8B-Instruct \\ \hline
\end{tabular}
\label{tab:verified_routes}
\end{table}

The change in routing settings for the same prompt resulted in a change in the model selection. This shows that the decision interface is affected by the selected policy and not just returning one single fixed model for those settings. The options cannot be taken to mean that GPT-5 is the best performing model, or Balanced routing is superior in terms of quality or costs to Legacy. These depend on generated responses and scores. What can be defended here is the difference in selection decisions resulting from those settings.

Both the privacy-related reactions involved choosing a model from the approved two entry models. Indeed, the ordinary reaction involved a change from GPT-4.1 in the balanced mode to an approved model in the private mode. It thus provides support to the limited finding that the choice of the private mode influenced the eligibility of the ordinary reaction. The sensitive balanced reaction was also in line with the restriction set up. However, the prompt as well as the flag used were different from the ordinary balanced reaction. It is impossible to determine the influence of the keyword in relation to the flag since the experiment cannot separate the two.

\subsection{Response to a Configured Sensitivity Keyword}
The third experiment sent the synthetic sensitive prompt directly to \texttt{/api/sensitivity}. The response labelled it sensitive and identified the word ``confidential'' as the matching keyword. The same word appeared in the recorded configuration. Table~\ref{tab:screening_experiment} presents the returned fields used in the analysis.

\begin{table}[!htbp]
\centering\small
\setstretch{1}\renewcommand{\arraystretch}{1.15}
\caption{Recorded response for the synthetic sensitivity example}
\begin{tabular}{|P{0.24\textwidth}|P{0.27\textwidth}|P{0.37\textwidth}|}\hline
\textbf{Response Field} & \textbf{Recorded Value} & \textbf{Interpretation} \\ \hline
HTTP status & 200 & The endpoint returned a successful response. \\ \hline
Sensitivity label & Sensitive & The supplied sample triggered a sensitivity rule. \\ \hline
Explanation & Keyword ``confidential'' matched. & The response identified a configured rule as the reason for the label. \\ \hline
\end{tabular}
\label{tab:screening_experiment}
\end{table}

Such a case indicates the existence of an established relationship between the specified keyword configuration and the classified content. Furthermore, an explanation allows tracing of the decision back to the applicable rule, which could help analyze why the given request was classified as sensitive. Such an outcome does not prove general detection capabilities. A sensitive request may include no keyword from the list, and regular text can use the concept of confidentiality without mentioning any confidential information. As there is only one confirmed positive case, we cannot measure precision and recall.

\section{Discussion}
The overall experiments have created a small but observable piece of the entire request flow as proposed. The prototype has provided the model information, has given policy-based selections and a rule-based justification of the sensitivity example. The configuration and response logs ensure that each observation is tested against the actual input used. This is an important contribution towards the FYP-I decision interface and behavior to maintain when integrating the selector with the gateway. However, the results have not demonstrated the full value of multi-objective advantage of the AetherFlow since a different model selected is not necessarily the better answer, cheaper or faster.

The API responses recorded also have predictions of quality and estimates of resource value. These have not been independently verified from the training data, benchmark entries, or actual usage of the providers, and hence are not included in the tables of results of the experiments. There was no scoring of any generated answers, nor was there a determination of cost and time for inference and generation. The difference is that while request fulfillment and explanation can be right, even if the prediction and/or response are wrong.

\section{Limitations and Further Evaluation}
This evidence is confined to a carefully curated set of prompts, where there was only one execution per prompt and setup. It does not cover variability between executions, the ability to work concurrently, behavior on failures of the providers, or performance on unseen task distributions. Additionally, the lack of the identifier of the server build used makes it more difficult to match the future setup to the one observed. The privacy concerns relate to the model names that are returned and do not cover the internal network traffic, local executions, or deletion of temporary files.

The FYP-II approach would make use of a tested dataset with well-defined answer labels, as well as a training, validation, and test split. The routing policies would be evaluated against baselines using the fixed models through the same query and generation settings. There would be scores for the generated answers, recording of the token usage and associated cost, and measurement of the response time. The privacy and cache evaluation would evaluate the approved destinations, data transfer, user isolation, deletion of the temporary mapping, and cache misdirection. Test cases and metrics for the integrated gateway would be designed.

\section{Conclusion}
The initial experiments show the behavior of the prototype in an observable way. The model configuration can be obtained; routing configurations will influence the selection; and sensitivity is the keyword provided, which leads to a classified outcome. Both privacy related selections that have been tested fall into the approved category. There are still issues related to answer quality, resource usage, and privacy left for evaluation in FYP-II.

\chapter{Conclusion}
\label{ch:conclusions}

AetherFlow solves the problem of model selection with regard to answer quality, costs, response time and privacy. This problem was studied using literature review, requirement analysis, proposed methodology, system design, prototype implementation and preliminary experiments. The literature demonstrated techniques for learned routing, semantic caching and efficient model serving to support the proposed gateway, and defined boundaries that need to be assessed in realistic settings. The requirements set out norms of conduct and specified that before routing preferences were used, privacy restrictions must be applied to determine which destinations would be allowed. The methodology and architecture explained how the client and other components such as privacy screening, cache, quality predictor, model adapters and storage services would work together. In FYPI, a prototype of the initial client and routing came up to investigate these decisions. The saved responses to the API calls indicated that the routing configurations influenced the models returned, and that the privacy-related configurations examined were part of the allowed set of configurations. There was also a synthetic sensitive prompt which yielded a classification related to a known keyword. These results are proof of prototype behavior, not proof of improvements in answer accuracy, cost, speed or comprehensive privacy protection. The project thus progressed in design and first implementation and the entire integrated scope has not yet been completed. The next step is to integrate the VS Code client, live model execution, secure key management, privacy services, and semantic cache. FYP-II will test the entire workflow on the checked datasets and record the usage, repeated measurements, and systematic testing, and present the evidence of the practical benefits and limitations of the approach proposed.

\section{Challenges and Work Left Out}
The problem is making sure that model selection leads to execution and results. The full log of the training and testing process has not yet been independently verified, which makes it impossible to draw any conclusions about the quality of predictions and performance on queries not seen before. The comparative analysis needs consistent scores of the models' responses, usage statistics, and consistent measurements. The routing results of the prototype cannot be used to confirm secure local processing or cache re-use.

\begin{table}[!htbp]
\centering\small
\setstretch{1}\small\setlength{\parskip}{0pt}\setlength{\tabcolsep}{5pt}\setlength{\extrarowheight}{1.5pt}\renewcommand{\arraystretch}{1.25}
\caption{Work deferred to FYP-II}
\begin{tabular}{|p{0.34\textwidth}|p{0.56\textwidth}|}\hline
\textbf{Item} & \textbf{Reason for Deferral} \\ \hline
Live local and cloud inference & Needs model adapters, credential protection and request flow completion. \\ \hline
Masking and Privacy Vault implementation & Needs verification of data transfer, data restoration and deletion mapping. \\ \hline
Two-step semantic cache & Needs privacy rules and isolation, and checking incorrect answer recycling. \\ \hline
Authentication, usage logging and recovery & Needs to depend on the connected gateway and storage services. \\ \hline
Performance evaluation and comparison & Needs verified data, generated answers, usage logs and timings. \\ \hline
Test cases and metrics development & To be designed and used for component testing and request flow. \\ \hline
\end{tabular}
\label{tab:deferred_work}
\end{table}

\section{Plan for FYP-II}
In FYP-II, the extension for VS Code will be coupled with the FastAPI gateway for handling chat and selected-code requests. Local Qwen and Llama models as well as the approved cloud-based models will be connected via different adapters. The user provided API keys and system managed API keys will be secured. The integration of privacy screening, Privacy Vault, and Qdrant caching will take place, while the request details will be recorded in PostgreSQL without including the sensitive information. The dataset will be reviewed by the team, quality predictor will be trained, and routing configuration will be tuned using the validation dataset. The testing process will include integration testing, privacy testing, cache correctness testing, timeout testing, and recovery testing.

\section{Suggestions for Future Work}
\begin{enumerate}
\setlength{\itemsep}{0pt}
\setlength{\parsep}{0pt}
\item Use uncertainty of prediction made by the router as a metric for switching to a more powerful model.
\item Perform model selection for various tasks, including queries where a less expensive model is correct and the most powerful model is wrong.
\item Detect sensitive data based on something other than keywords and test the cache against similar looking queries with different answers.
\item Investigate and fix slow responses of the local models.
\item Re-test updated models with realistic requests.
\end{enumerate}

The report establishes AetherFlow's requirements, literature-based approach, system design, and initial prototype. FYP-II will complete the gateway and evaluate its practical benefits through reproducible experiments, addressing the integration and evidence gaps identified throughout this report.
